\documentclass[preview]{standalone}

\usepackage{amsmath}
\usepackage{amssymb}
\usepackage{stellar}
\usepackage{definitions}
\usepackage{bettelini}

\begin{document}

\id{linearalgebra-quadratic-forms}
\genpage

\section{Quadratic forms}

\begin{snippetdefinition}{symmetric-bilinear-form-definition}{Symmetric bilinear form}
    Let \(V\) be a real \vectorspace. A \bilinearmap
    \[
        B\colon V\cartesianprod V\fromto\realnumbers
    \]
    is a \emph{symmetric bilinear form} if
    \[
        B(u,v)=B(v,u)
    \]
    for every \(u,v\in V\).
\end{snippetdefinition}

\begin{snippetdefinition}{quadratic-form-definition}{Quadratic form}
    Let \(V\) be a real \vectorspace. A \emph{quadratic form} on \(V\) is a
    \function
    \[
        Q\colon V\fromto\realnumbers
    \]
    such that:
    \begin{enumerate}
        \item \(Q(\lambda v)=\lambda^2Q(v)\) for every \(v\in V\) and
            every \(\lambda\in\realnumbers\);
        \item the map
            \[
                B_Q\colon V\cartesianprod V\fromto\realnumbers,
                \qquad
                B_Q(u,v)=\frac{Q(u+v)-Q(u)-Q(v)}{2}
            \]
            is bilinear.
    \end{enumerate}
    The bilinear map \(B_Q\) is called the
    \emph{associated bilinear form} of \(Q\).
\end{snippetdefinition}

\begin{snippetproposition}{associated-bilinear-form-is-symmetric}{The associated bilinear form is symmetric}
    Let \(V\) be a real \vectorspace, and let \(Q\colon V\fromto\realnumbers\)
    be a \quadraticform with \associatedbilinearform \(B_Q\). Then \(B_Q\) is
    a \symmetricbilinearform.
\end{snippetproposition}

\begin{snippetproof}{associated-bilinear-form-is-symmetric-proof}{associated-bilinear-form-is-symmetric}{The associated bilinear form is symmetric}
    For every \(u,v\in V\),
    \[
        B_Q(u,v)
        =
        \frac{Q(u+v)-Q(u)-Q(v)}{2}
        =
        \frac{Q(v+u)-Q(v)-Q(u)}{2}
        =
        B_Q(v,u).
    \]
\end{snippetproof}

\begin{snippetexample}{inner-product-square-norm-quadratic-form}{Squared norm as a quadratic form}
    Let \((V,\langle\cdot,\cdot\rangle)\) be a real \innerproductspace. The
    map
    \[
        Q\colon V\fromto\realnumbers,
        \qquad
        Q(v)=\|v\|^2
    \]
    is a \quadraticform. Its \associatedbilinearform is the
    \innerproduct:
    \[
        B_Q(u,v)=\langle u,v\rangle.
    \]
    Indeed,
    \[
        Q(\lambda v)=\|\lambda v\|^2=\lambda^2\|v\|^2=\lambda^2Q(v),
    \]
    and
    \[
        \frac{\|u+v\|^2-\|u\|^2-\|v\|^2}{2}
        =
        \langle u,v\rangle.
    \]
\end{snippetexample}

\begin{snippetexample}{spacetime-type-quadratic-form-example}{A spacetime-type quadratic form}
    Fix \(c\in\realnumbers\). The map
    \[
        Q\colon\realnumbers^4\fromto\realnumbers,
        \qquad
        Q(x,y,z,t)=c^2t^2-x^2-y^2-z^2
    \]
    is a \quadraticform. Its \associatedbilinearform is
    \[
        B_Q\bigl((x,y,z,t),(x',y',z',t')\bigr)
        =
        c^2tt'-xx'-yy'-zz'.
    \]
    This form takes positive, negative, and zero values on nonzero vectors.
\end{snippetexample}

\begin{snippetexample}{polynomial-evaluation-quadratic-form-example}{Quadratic form from polynomial evaluations}
    Let \(V=\mathcal{P}_3(\realnumbers)\). The map
    \[
        Q\colon V\fromto\realnumbers,
        \qquad
        Q(p)=p(0)^2+p(1)^2+p(2)^2
    \]
    is a \quadraticform. Its \associatedbilinearform is
    \[
        B_Q(p,q)=p(0)q(0)+p(1)q(1)+p(2)q(2).
    \]
    However, \(Q\) is not \positivedefiniteform[positive definite], because
    the nonzero polynomial
    \[
        p(x)=x(x-1)(x-2)
    \]
    satisfies \(Q(p)=0\).
\end{snippetexample}

\begin{snippetexample}{symmetric-matrix-quadratic-form-example}{Quadratic form associated to a symmetric matrix}
    Let \(S\in\matrices_{n\times n}(\realnumbers)\) be a
    \symmetricmatrix. Define
    \[
        Q_S\colon\realnumbers^n\fromto\realnumbers,
        \qquad
        Q_S(x)=x^\transpose Sx.
    \]
    Then \(Q_S\) is a \quadraticform. Its \associatedbilinearform is
    \[
        B_{Q_S}(x,y)=x^\transpose Sy.
    \]
\end{snippetexample}

\begin{snippetproof}{symmetric-matrix-quadratic-form-example-proof}{symmetric-matrix-quadratic-form-example}{Quadratic form associated to a symmetric matrix}
    Since \(S\) is symmetric,
    \[
        Q_S(x+y)-Q_S(x)-Q_S(y)
        =
        x^\transpose Sy+y^\transpose Sx
        =
        2x^\transpose Sy.
    \]
    Hence
    \[
        B_{Q_S}(x,y)
        =
        \frac{Q_S(x+y)-Q_S(x)-Q_S(y)}{2}
        =
        x^\transpose Sy.
    \]
\end{snippetproof}

\begin{snippetexample}{trace-square-quadratic-form-example}{Trace of the square}
    On \(\matrices_{n\times n}(\realnumbers)\), define
    \[
        Q(A)=\trace(A^2).
    \]
    Then \(Q\) is a \quadraticform. Its \associatedbilinearform is
    \[
        B_Q(A,C)=\trace(AC).
    \]
    Indeed,
    \[
        Q(A+C)-Q(A)-Q(C)
        =
        \trace(AC+CA)
        =
        2\trace(AC).
    \]
\end{snippetexample}

\section{Sign}

\begin{snippetdefinition}{quadratic-form-sign-definition}{Sign of a quadratic form}
    Let \(V\) be a real \vectorspace, and let
    \(Q\colon V\fromto\realnumbers\) be a \quadraticform.
    \begin{enumerate}
        \item \(Q\) is \emph{positive definite} if \(Q(v)>0\) for every
            nonzero \(v\in V\).
        \item \(Q\) is \emph{positive semidefinite} if \(Q(v)\geq 0\) for
            every \(v\in V\), and there exists a nonzero \(v\in V\) such
            that \(Q(v)=0\).
        \item \(Q\) is \emph{indefinite} if there exist \(v,w\in V\) such
            that \(Q(v)<0<Q(w)\).
        \item \(Q\) is \emph{negative semidefinite} if \(Q(v)\leq 0\) for
            every \(v\in V\), and there exists a nonzero \(v\in V\) such
            that \(Q(v)=0\).
        \item \(Q\) is \emph{negative definite} if \(Q(v)<0\) for every
            nonzero \(v\in V\).
    \end{enumerate}
\end{snippetdefinition}

\begin{snippetexample}{quadratic-form-sign-examples}{Signs of the previous examples}
    The \quadraticform \(Q(v)=\|v\|^2\) associated to a real
    \innerproduct is \positivedefiniteform[positive definite].

    The \quadraticform
    \[
        Q(x,y,z,t)=c^2t^2-x^2-y^2-z^2
    \]
    is \indefiniteform[indefinite].

    The \quadraticform
    \[
        Q(p)=p(0)^2+p(1)^2+p(2)^2
    \]
    on \(\mathcal{P}_3(\realnumbers)\) is
    \positivesemidefiniteform[positive semidefinite].

    The \quadraticform \(Q(A)=\trace(A^2)\) on
    \(\matrices_{n\times n}(\realnumbers)\) is
    \indefiniteform[indefinite] when \(n\geq 2\). For example, on
    \(\matrices_{2\times 2}(\realnumbers)\),
    \[
        Q(\identmatrix{2})=2>0,
    \]
    while, for
    \[
        R=
        \begin{pmatrix}
            0 & -1\\
            1 & 0
        \end{pmatrix},
    \]
    we have \(R^2=-\identmatrix{2}\), and therefore \(Q(R)=-2<0\).
\end{snippetexample}

\section{Bilinear correspondence}

\begin{snippetproposition}{quadratic-form-bilinear-correspondence}{Quadratic forms and symmetric bilinear forms}
    Let \(V\) be a real \vectorspace. If \(Q\colon V\fromto\realnumbers\) is
    a \quadraticform and \(B_Q\) is its \associatedbilinearform, then
    \[
        Q(v)=B_Q(v,v)
    \]
    for every \(v\in V\).

    Conversely, if \(B\colon V\cartesianprod V\fromto\realnumbers\) is a
    \symmetricbilinearform, then
    \[
        Q_B(v)=B(v,v)
    \]
    defines a \quadraticform on \(V\), and the
    \associatedbilinearform of \(Q_B\) is \(B\).
\end{snippetproposition}

\begin{snippetproof}{quadratic-form-bilinear-correspondence-proof}{quadratic-form-bilinear-correspondence}{Quadratic forms and symmetric bilinear forms}
    If \(Q\) is a \quadraticform, then
    \[
        B_Q(v,v)
        =
        \frac{Q(2v)-Q(v)-Q(v)}{2}
        =
        \frac{4Q(v)-2Q(v)}{2}
        =
        Q(v).
    \]

    Conversely, let \(B\) be a \symmetricbilinearform and define
    \(Q_B(v)=B(v,v)\). Then
    \[
        Q_B(\lambda v)
        =
        B(\lambda v,\lambda v)
        =
        \lambda^2B(v,v)
        =
        \lambda^2Q_B(v).
    \]
    The \associatedbilinearform of \(Q_B\) is
    \[
        \widetilde{B}(u,v)
        =
        \frac{Q_B(u+v)-Q_B(u)-Q_B(v)}{2}.
    \]
    Since \(B\) is bilinear and symmetric,
    \[
        \widetilde{B}(u,v)
        =
        \frac{B(u,v)+B(v,u)}{2}
        =
        B(u,v).
    \]
    Hence \(\widetilde{B}=B\).
\end{snippetproof}

\section{Associated matrix}

\begin{snippetdefinition}{quadratic-form-associated-matrix-definition}{Matrix associated to a quadratic form}
    Let \(V\) be a finite-dimensional real \vectorspace, let
    \(Q\colon V\fromto\realnumbers\) be a \quadraticform, and let \(B_Q\) be
    the \associatedbilinearform of \(Q\). If
    \[
        \mathcal{B}=\{v_1,\ldots,v_n\}
    \]
    is a \basis of \(V\), the \emph{matrix associated to \(Q\) with respect
    to \(\mathcal{B}\)} is
    \[
        S_{\mathcal{B}}(Q)=(s_{ij})\in\matrices_{n\times n}(\realnumbers),
        \qquad
        s_{ij}=B_Q(v_i,v_j).
    \]
\end{snippetdefinition}

\begin{snippetproposition}{quadratic-form-associated-matrix-is-symmetric}{Associated matrices of real quadratic forms are symmetric}
    Let \(V\) be a finite-dimensional real \vectorspace, let
    \(Q\colon V\fromto\realnumbers\) be a \quadraticform, and let
    \(\mathcal{B}\) be a \basis of \(V\). Then the
    \quadraticformmatrix of \(Q\) with respect to \(\mathcal{B}\) is a
    \symmetricmatrix.
\end{snippetproposition}

\begin{snippetproof}{quadratic-form-associated-matrix-is-symmetric-proof}{quadratic-form-associated-matrix-is-symmetric}{Associated matrices of real quadratic forms are symmetric}
    If \(S_{\mathcal{B}}(Q)=(s_{ij})\), then
    \[
        s_{ij}=B_Q(v_i,v_j)=B_Q(v_j,v_i)=s_{ji}
    \]
    for every \(i,j\). Therefore
    \[
        S_{\mathcal{B}}(Q)^\transpose=S_{\mathcal{B}}(Q).
    \]
\end{snippetproof}

\begin{snippetproposition}{quadratic-form-coordinate-formula}{Coordinate formula for a quadratic form}
    Let \(V\) be a finite-dimensional real \vectorspace, let
    \(Q\colon V\fromto\realnumbers\) be a \quadraticform, let
    \(\mathcal{B}=\{v_1,\ldots,v_n\}\) be a \basis of \(V\), and let
    \(S_{\mathcal{B}}(Q)\) be the \quadraticformmatrix of \(Q\). If
    \[
        v=x_1v_1+\cdots+x_nv_n
    \]
    and
    \[
        [v]_{\mathcal{B}}=
        \begin{pmatrix}
            x_1\\
            \vdots\\
            x_n
        \end{pmatrix},
    \]
    then
    \[
        Q(v)
        =
        [v]_{\mathcal{B}}^\transpose
        S_{\mathcal{B}}(Q)
        [v]_{\mathcal{B}}.
    \]
\end{snippetproposition}

\begin{snippetproof}{quadratic-form-coordinate-formula-proof}{quadratic-form-coordinate-formula}{Coordinate formula for a quadratic form}
    Since \(Q(v)=B_Q(v,v)\) and \(B_Q\) is bilinear,
    \[
        Q(v)
        =
        B_Q\left(\sum_{i=1}^n x_iv_i,\sum_{j=1}^n x_jv_j\right)
        =
        \sum_{i,j=1}^n x_ix_jB_Q(v_i,v_j).
    \]
    This is exactly the \matrix product
    \[
        [v]_{\mathcal{B}}^\transpose
        S_{\mathcal{B}}(Q)
        [v]_{\mathcal{B}}.
    \]
\end{snippetproof}

\section{Congruence}

\includesnpt{similar-matrices-definition}
\includesnpt{congruent-matrices-definition}

\begin{snippetnote}{similarity-and-congruence-are-different}{Similarity and congruence are different}
    \similarmatrices[Similarity] and \congruentmatrices[congruence] are
    different relations. The identity \matrix is similar only to itself,
    because
    \[
        P^{-1}\identmatrix{n}P=\identmatrix{n}
    \]
    for every invertible \(P\in\matrices_{n\times n}(\realnumbers)\).

    On the other hand, many \matrix[matrices] different from \(\identmatrix{n}\) are
    congruent to \(\identmatrix{n}\). For example, if
    \[
        P=
        \begin{pmatrix}
            2 & 0\\
            0 & 3
        \end{pmatrix},
    \]
    then
    \[
        P^\transpose\identmatrix{2}P
        =
        \begin{pmatrix}
            4 & 0\\
            0 & 9
        \end{pmatrix}.
    \]
\end{snippetnote}

\begin{snippetdefinition}{congruent-quadratic-forms-definition}{Congruent quadratic forms}
    Let \(V\) be a finite-dimensional real \vectorspace. Two
    \quadraticform[quadratic forms]
    \[
        Q,Q'\colon V\fromto\realnumbers
    \]
    are \emph{congruent} if there exist two \basis[bases]
    \(\mathcal{B}\) and \(\mathcal{B}'\) of \(V\) such that the
    \quadraticformmatrix[associated matrices]
    \[
        S_{\mathcal{B}}(Q),
        \qquad
        S_{\mathcal{B}'}(Q')
    \]
    are \congruentmatrices.
\end{snippetdefinition}

\begin{snippet}{quadratic-forms-congruence-goals}
    Congruence is the natural relation for studying real
    \quadraticform[quadratic forms]: changing coordinates changes the
    associated \matrix by congruence. The main goals are to determine the
    \quadraticformsign of a \quadraticform and to classify
    \quadraticform[quadratic forms] up to congruence.
\end{snippet}

\section{Diagonalization}

\begin{snippetproposition}{quadratic-form-diagonal-form}{Diagonal form of a real quadratic form}
    Let \(V\) be a finite-dimensional real \vectorspace, and let
    \(Q\colon V\fromto\realnumbers\) be a \quadraticform. There exists a
    \basis
    \[
        \mathcal{E}=\{e_1,\ldots,e_n\}
    \]
    of \(V\) such that, for every
    \[
        v=a_1e_1+\cdots+a_ne_n,
    \]
    one has
    \[
        Q(v)=\lambda_1a_1^2+\cdots+\lambda_na_n^2,
    \]
    where \(\lambda_1,\ldots,\lambda_n\) are the
    \eigenvalue[eigenvalues] of a \quadraticformmatrix of \(Q\).
\end{snippetproposition}

\begin{snippetproof}{quadratic-form-diagonal-form-proof}{quadratic-form-diagonal-form}{Diagonal form of a real quadratic form}
    Let \(\mathcal{B}\) be a \basis of \(V\), and let
    \(S=S_{\mathcal{B}}(Q)\). The \matrix \(S\) is symmetric. By the real
    \spectraltheorem, there exists an \orthogonalmatrix \(P\) such that
    \[
        P^\transpose SP=D,
    \]
    where \(D=\operatorname{diag}(\lambda_1,\ldots,\lambda_n)\).

    Interpret \(P\) as a change-of-basis \matrix. In the resulting
    \basis \(\mathcal{E}\), the \quadraticformmatrix of \(Q\) is \(D\).
    Therefore, if
    \[
        [v]_{\mathcal{E}}=
        \begin{pmatrix}
            a_1\\
            \vdots\\
            a_n
        \end{pmatrix},
    \]
    then
    \[
        Q(v)
        =
        [v]_{\mathcal{E}}^\transpose
        D
        [v]_{\mathcal{E}}
        =
        \lambda_1a_1^2+\cdots+\lambda_na_n^2.
    \]
\end{snippetproof}

\begin{snippettheorem}{sylvester-normal-form}{Sylvester normal form}
    Let \(V\) be a finite-dimensional real \vectorspace, and let
    \(Q\colon V\fromto\realnumbers\) be a \quadraticform. There exists a
    \basis of \(V\) in which
    \[
        Q(x_1,\ldots,x_n)
        =
        x_1^2+\cdots+x_p^2
        -
        x_{p+1}^2-\cdots-x_{p+q}^2,
    \]
    where \(p,q\geq 0\). The remaining coordinates, if any, have
    coefficient \(0\). Equivalently, in a suitable \basis, the
    \quadraticformmatrix of \(Q\) is
    \[
        \begin{pmatrix}
            \identmatrix{p} & 0 & 0\\
            0 & -\identmatrix{q} & 0\\
            0 & 0 & 0
        \end{pmatrix}.
    \]
\end{snippettheorem}

\begin{snippetproof}{sylvester-normal-form-proof}{sylvester-normal-form}{Sylvester normal form}
    By the diagonal form proposition, there is a \basis in which
    \[
        Q(x_1,\ldots,x_n)=\lambda_1x_1^2+\cdots+\lambda_nx_n^2.
    \]
    Reorder the \basis so that
    \[
        \lambda_1,\ldots,\lambda_p>0,
        \qquad
        \lambda_{p+1},\ldots,\lambda_{p+q}<0,
    \]
    and
    \[
        \lambda_{p+q+1}=\cdots=\lambda_n=0.
    \]
    For each positive \(\lambda_i\), replace \(e_i\) by
    \(e_i/\sqrt{\lambda_i}\). For each negative \(\lambda_i\), replace
    \(e_i\) by \(e_i/\sqrt{-\lambda_i}\). The positive coefficients become
    \(1\), the negative coefficients become \(-1\), and the zero
    coefficients remain zero.
\end{snippetproof}

\section{Sylvester theorem}

\begin{snippetdefinition}{positive-negative-subspaces-definition}{Positive and negative subspaces}
    Let \(V\) be a real \vectorspace, and let
    \(Q\colon V\fromto\realnumbers\) be a \quadraticform. A subspace
    \(U\subseteq V\) is \emph{positive for \(Q\)} if the restriction
    \(Q|_U\) is \positivedefiniteform[positive definite], that is,
    \[
        Q(u)>0
    \]
    for every nonzero \(u\in U\). It is \emph{negative for \(Q\)} if
    \(Q|_U\) is \negativedefiniteform[negative definite], that is,
    \[
        Q(u)<0
    \]
    for every nonzero \(u\in U\).
\end{snippetdefinition}

\begin{snippetdefinition}{quadratic-form-maximal-positive-negative-dimensions}{Maximal positive and negative dimensions}
    Let \(V\) be a finite-dimensional real \vectorspace, and let
    \(Q\colon V\fromto\realnumbers\) be a \quadraticform. Define
    \[
        m_+(Q)
        =
        \max\{\lineardim U \suchthat U\subseteq V,\ Q|_U
        \text{ is positive definite}\},
    \]
    and
    \[
        m_-(Q)
        =
        \max\{\lineardim U \suchthat U\subseteq V,\ Q|_U
        \text{ is negative definite}\}.
    \]
\end{snippetdefinition}

\begin{snippettheorem}{sylvester-law-of-inertia}{Sylvester's law of inertia}
    Let \(V\) be a finite-dimensional real \vectorspace, and let
    \(Q\colon V\fromto\realnumbers\) be a \quadraticform. Suppose that, in
    some \basis, \(Q\) is represented by the diagonal \matrix
    \[
        \operatorname{diag}(\lambda_1,\ldots,\lambda_n).
    \]
    Let
    \[
        p=\#\{i \suchthat \lambda_i>0\},
        \qquad
        q=\#\{i \suchthat \lambda_i<0\},
        \qquad
        r=\#\{i \suchthat \lambda_i=0\}.
    \]
    Then \(p,q,r\) do not depend on the chosen \basis. Equivalently, two
    congruent real \symmetricmatrix[symmetric matrices] have the same numbers
    of positive, negative, and zero \eigenvalue[eigenvalues].
\end{snippettheorem}

\begin{snippetproof}{sylvester-law-of-inertia-proof}{sylvester-law-of-inertia}{Sylvester's law of inertia}
    We prove that the number \(p\) of positive coefficients is intrinsic.
    The proof for \(q\) is the same applied to \(-Q\), and then
    \(r=n-p-q\).

    Work in a diagonal \basis and reorder it so that
    \[
        Q(x_1,\ldots,x_n)
        =
        \lambda_1x_1^2+\cdots+\lambda_nx_n^2,
    \]
    with
    \[
        \lambda_1,\ldots,\lambda_p>0,
        \qquad
        \lambda_{p+1},\ldots,\lambda_n\leq 0.
    \]
    Let
    \[
        U_+
        =
        \{(x_1,\ldots,x_n) \suchthat x_{p+1}=\cdots=x_n=0\}.
    \]
    Then \(\lineardim U_+=p\), and \(Q\) is positive definite on \(U_+\).
    Hence \(m_+(Q)\geq p\).

    Conversely, let \(U\subseteq V\) be a subspace with \(\lineardim U>p\).
    Consider
    \[
        \pi_+\colon U\fromto\realnumbers^p,
        \qquad
        \pi_+(x_1,\ldots,x_n)=(x_1,\ldots,x_p).
    \]
    Since \(\lineardim U>p=\lineardim\realnumbers^p\), the map \(\pi_+\)
    is not injective. Hence there exists a nonzero \(u\in U\) such that
    \(\pi_+(u)=0\). Thus
    \[
        u=(0,\ldots,0,x_{p+1},\ldots,x_n),
    \]
    and
    \[
        Q(u)
        =
        \lambda_{p+1}x_{p+1}^2+\cdots+\lambda_nx_n^2
        \leq 0.
    \]
    Therefore \(Q\) is not positive definite on \(U\). Every positive
    subspace has dimension at most \(p\), so \(m_+(Q)\leq p\).

    Hence \(m_+(Q)=p\). Since \(m_+(Q)\) is defined intrinsically from
    \(Q\), the number \(p\) is independent of the diagonal \basis. Applying
    the same argument to \(-Q\) proves the invariance of \(q\), and then
    \(r=n-p-q\) is invariant as well.
\end{snippetproof}

\begin{snippetcorollary}{quadratic-forms-congruence-classification}{Classification by congruence}
    Every finite-dimensional real \quadraticform is congruent to a unique
    form of the type
    \[
        x_1^2+\cdots+x_p^2
        -
        x_{p+1}^2-\cdots-x_{p+q}^2,
    \]
    up to reordering coordinates. Equivalently, every real
    \symmetricmatrix is congruent to a unique \matrix of the form
    \[
        \begin{pmatrix}
            \identmatrix{p} & 0 & 0\\
            0 & -\identmatrix{q} & 0\\
            0 & 0 & 0_{r\times r}
        \end{pmatrix},
    \]
    where \(p,q,r\) are determined by the \quadraticform.
\end{snippetcorollary}

\begin{snippetproof}{quadratic-forms-congruence-classification-proof}{quadratic-forms-congruence-classification}{Classification by congruence}
    Existence follows from the \sylvesternormalform. Uniqueness follows from
    the \sylvesterlaw: \(p\) is the maximal dimension of a positive
    subspace, \(q\) is the maximal dimension of a negative subspace, and
    \(r=n-p-q\).
\end{snippetproof}

\begin{snippetdefinition}{inertia-indices-definition}{Indices of inertia}
    Let \(V\) be a finite-dimensional real \vectorspace, and let
    \(Q\colon V\fromto\realnumbers\) be a \quadraticform. If the
    \sylvesternormalform of \(Q\) is
    \[
        x_1^2+\cdots+x_p^2
        -
        x_{p+1}^2-\cdots-x_{p+q}^2,
    \]
    with \(r\) remaining zero coefficients, then the triple
    \[
        (p,q,r)
    \]
    is called the \emph{signature} or the \emph{indices of inertia} of
    \(Q\). The number \(p\) is the positive index, \(q\) is the negative
    index, and \(r\) is the null index.
\end{snippetdefinition}

\section{Matrix representations}

\begin{snippetexample}{polynomial-quadratic-form-associated-matrix-example}{Quadratic form on polynomials}
    Let \(V=\mathcal{P}_2(\realnumbers)\), and define
    \[
        Q\colon V\fromto\realnumbers,
        \qquad
        Q(p)=p(0)p(1)+p(-1)p(2).
    \]
    The \associatedbilinearform of \(Q\) is
    \[
        B_Q(p,q)
        =
        \frac{
            p(0)q(1)+q(0)p(1)+p(-1)q(2)+p(2)q(-1)
        }{2}.
    \]
    With respect to the \basis
    \[
        \mathcal{B}=\{1,x,x^2\},
    \]
    the \quadraticformmatrix of \(Q\) is
    \[
        S_{\mathcal{B}}(Q)
        =
        \begin{pmatrix}
            2 & 1 & 3\\
            1 & -2 & -1\\
            3 & -1 & 4
        \end{pmatrix}.
    \]
\end{snippetexample}

\begin{snippetproof}{polynomial-quadratic-form-associated-matrix-example-proof}{polynomial-quadratic-form-associated-matrix-example}{Quadratic form on polynomials}
    If \(\mathcal{B}=\{v_1,v_2,v_3\}=\{1,x,x^2\}\), then
    \[
        S_{\mathcal{B}}(Q)=(B_Q(v_i,v_j))_{ij}.
    \]
    For instance,
    \[
        B_Q(1,1)=\frac{1+1+1+1}{2}=2,
    \]
    and
    \[
        B_Q(x^2,x)=\frac{-4+2}{2}=-1.
    \]
    Computing all entries gives the displayed symmetric \matrix.
\end{snippetproof}

\begin{snippetproposition}{quadratic-form-matrix-change-of-basis}{Change of basis for associated matrices}
    Let \(V\) be a finite-dimensional real \vectorspace, let
    \(Q\colon V\fromto\realnumbers\) be a \quadraticform, and let
    \(\mathcal{B}\) and \(\mathcal{E}\) be \basis[bases] of \(V\). If
    \[
        A=M(\identityfunc_V,\mathcal{E},\mathcal{B}),
    \]
    then
    \[
        S_{\mathcal{E}}(Q)
        =
        A^\transpose S_{\mathcal{B}}(Q)A.
    \]
\end{snippetproposition}

\begin{snippetproof}{quadratic-form-matrix-change-of-basis-proof}{quadratic-form-matrix-change-of-basis}{Change of basis for associated matrices}
    The \matrix \(A\) transforms coordinates from \(\mathcal{E}\) to
    \(\mathcal{B}\):
    \[
        [v]_{\mathcal{B}}=A[v]_{\mathcal{E}}.
    \]
    Therefore
    \begin{align*}
        Q(v)
        &=
        [v]_{\mathcal{B}}^\transpose
        S_{\mathcal{B}}(Q)
        [v]_{\mathcal{B}}\\
        &=
        (A[v]_{\mathcal{E}})^\transpose
        S_{\mathcal{B}}(Q)
        (A[v]_{\mathcal{E}})\\
        &=
        [v]_{\mathcal{E}}^\transpose
        A^\transpose S_{\mathcal{B}}(Q)A
        [v]_{\mathcal{E}}.
    \end{align*}
    By uniqueness of the \quadraticformmatrix in the \basis
    \(\mathcal{E}\), we get
    \[
        S_{\mathcal{E}}(Q)=A^\transpose S_{\mathcal{B}}(Q)A.
    \]
\end{snippetproof}

\begin{snippetproposition}{congruence-quadratic-form-interpretation}{Interpretation of congruence}
    Let \(B,C\in\matrices_{n\times n}(\realnumbers)\) be
    \symmetricmatrix[symmetric matrices]. Then \(B\) and \(C\) are
    \congruentmatrices if and only if they represent the same
    \quadraticform on \(\realnumbers^n\) with respect to two different
    \basis[bases].
\end{snippetproposition}

\begin{snippetproof}{congruence-quadratic-form-interpretation-proof}{congruence-quadratic-form-interpretation}{Interpretation of congruence}
    If \(B\) and \(C\) represent the same \quadraticform \(Q\) in two
    different \basis[bases], the change-of-basis formula gives an invertible
    \matrix \(A\) such that
    \[
        C=A^\transpose BA.
    \]
    Hence \(B\) and \(C\) are \congruentmatrices.

    Conversely, suppose that \(C=A^\transpose BA\) for an invertible
    \matrix \(A\). Consider
    \[
        Q_B(x)=x^\transpose Bx
    \]
    with respect to the canonical \basis of \(\realnumbers^n\). If the
    columns of \(A\) are used as a new \basis, then the
    \quadraticformmatrix of \(Q_B\) in that new \basis is
    \[
        A^\transpose BA=C.
    \]
\end{snippetproof}

\begin{snippetexample}{xy-quadratic-form-congruence-example}{The quadratic form 2xy}
    Consider
    \[
        Q\colon\realnumbers^2\fromto\realnumbers,
        \qquad
        Q(x,y)=2xy.
    \]
    Its \associatedbilinearform is
    \[
        B_Q((x,y),(z,w))=xw+yz.
    \]
    In the canonical \basis, its \quadraticformmatrix is
    \[
        \begin{pmatrix}
            0 & 1\\
            1 & 0
        \end{pmatrix}.
    \]
    With respect to the \basis
    \[
        \mathcal{B}
        =
        \left\{
            \left(\frac{1}{2},\frac{1}{2}\right),
            \left(\frac{1}{2},-\frac{1}{2}\right)
        \right\},
    \]
    the \quadraticformmatrix is
    \[
        \begin{pmatrix}
            1 & 0\\
            0 & -1
        \end{pmatrix}.
    \]
    With respect to the \basis
    \[
        \mathcal{C}=\{(1,1),(1,-1)\},
    \]
    the \quadraticformmatrix is
    \[
        \begin{pmatrix}
            2 & 0\\
            0 & -2
        \end{pmatrix}.
    \]
    These \matrix[matrices] are different, but they are all congruent because they
    represent the same \quadraticform in different \basis[bases].
\end{snippetexample}

\begin{snippetnote}{positive-diagonal-matrices-congruent-to-identity}{Positive diagonal matrices are congruent to the identity}
    Let
    \[
        D=
        \begin{pmatrix}
            \lambda_1 & & 0\\
            & \ddots & \\
            0 & & \lambda_n
        \end{pmatrix},
        \qquad
        \lambda_i>0.
    \]
    Then \(D\) is congruent to \(\identmatrix{n}\). Indeed, with
    \[
        P=
        \begin{pmatrix}
            \sqrt{\lambda_1} & & 0\\
            & \ddots & \\
            0 & & \sqrt{\lambda_n}
        \end{pmatrix},
    \]
    one has
    \[
        P^\transpose\identmatrix{n}P=D.
    \]
\end{snippetnote}

\section{Eigenvalue criterion}

\begin{snippetproposition}{quadratic-form-sign-eigenvalue-criterion}{Sign criterion through eigenvalues}
    Let \(V\) be a finite-dimensional real \vectorspace, let
    \(Q\colon V\fromto\realnumbers\) be a \quadraticform, and let \(S\) be
    any \quadraticformmatrix of \(Q\). Then:
    \begin{enumerate}
        \item \(Q\) is \positivedefiniteform[positive definite] if and only
            if every \eigenvalue of \(S\) is positive;
        \item \(Q\) is \positivesemidefiniteform[positive semidefinite] if
            and only if every \eigenvalue of \(S\) is non-negative and at
            least one is zero;
        \item \(Q\) is \negativedefiniteform[negative definite] if and only
            if every \eigenvalue of \(S\) is negative;
        \item \(Q\) is \negativesemidefiniteform[negative semidefinite] if
            and only if every \eigenvalue of \(S\) is non-positive and at
            least one is zero;
        \item \(Q\) is \indefiniteform[indefinite] if and only if \(S\) has
            at least one positive \eigenvalue and at least one negative
            \eigenvalue.
    \end{enumerate}
\end{snippetproposition}

\begin{snippetproof}{quadratic-form-sign-eigenvalue-criterion-proof}{quadratic-form-sign-eigenvalue-criterion}{Sign criterion through eigenvalues}
    By the diagonal form proposition, there is a \basis in which
    \[
        Q(x_1,\ldots,x_n)
        =
        \lambda_1x_1^2+\cdots+\lambda_nx_n^2,
    \]
    where \(\lambda_1,\ldots,\lambda_n\) are the \eigenvalue[eigenvalues]
    of \(S\).
    The sign of \(Q\) is therefore read from the signs of the coefficients
    \(\lambda_i\). For example, if every \(\lambda_i>0\), then every
    nonzero vector has \(Q(v)>0\). Conversely, if some \(\lambda_i\leq 0\),
    the corresponding \basis vector gives a nonzero \(v\) with
    \(Q(v)\leq 0\). The other cases are analogous.
\end{snippetproof}

\section{Exercises}

\begin{snippetexercise}{quadratic-form-parameter-congruence-exercise}{Congruence of two parameter-dependent quadratic forms}
    Consider the \quadraticform[quadratic forms]
    \[
        Q,Q_k\colon\realnumbers^3\fromto\realnumbers
    \]
    defined by
    \[
        Q(x,y,z)=5x^2-y^2+z^2+4xy+6xz
    \]
    and
    \[
        Q_k(x,y,z)=2(xy+kyz-kxz).
    \]
    Determine the values of \(k\in\realnumbers\) for which \(Q\) and
    \(Q_k\) are \congruentquadraticforms.
\end{snippetexercise}

\begin{snippetsolution}{quadratic-form-parameter-congruence-solution}{Congruence of two parameter-dependent quadratic forms}
    The \quadraticformmatrix of \(Q\) in the canonical \basis is
    \[
        S=
        \begin{pmatrix}
            5 & 2 & 3\\
            2 & -1 & 0\\
            3 & 0 & 1
        \end{pmatrix},
    \]
    because the mixed terms are split symmetrically:
    \(4xy=2\cdot 2xy\) and \(6xz=2\cdot 3xz\). Its
    \characteristicpolynomial is
    \[
        p_S(t)=\det(t\identmatrix{3}-S)=t(t-7)(t+2).
    \]
    Thus the \signature of \(Q\) is \((1,1,1)\).

    The \quadraticformmatrix of \(Q_k\) in the canonical \basis is
    \[
        T_k=
        \begin{pmatrix}
            0 & 1 & -k\\
            1 & 0 & k\\
            -k & k & 0
        \end{pmatrix}.
    \]
    Its \characteristicpolynomial is
    \[
        p_{T_k}(t)=\det(t\identmatrix{3}-T_k)=(t-1)(t^2+t-2k^2).
    \]
    For \(Q_k\) to have the same \signature as \(Q\), the \matrix \(T_k\)
    must have a zero \eigenvalue. This happens if and only if
    \(p_{T_k}(0)=0\), namely
    \[
        2k^2=0.
    \]
    Hence \(k=0\).

    If \(k=0\), then
    \[
        p_{T_0}(t)=t(t-1)(t+1),
    \]
    so the \signature is \((1,1,1)\). If \(k\neq 0\), then \(0\) is not an
    \eigenvalue and \(t^2+t-2k^2\) has one positive root and one negative
    root. Together with the \eigenvalue \(1\), this gives \signature
    \((2,1,0)\). By \sylvesterlaw,
    \[
        Q \text{ and } Q_k \text{ are congruent}
        \qquad\Longleftrightarrow\qquad
        k=0.
    \]
\end{snippetsolution}

\begin{snippetexercise}{symmetric-bilinear-form-inner-product-exercise}{A symmetric bilinear form on R3}
    Consider the bilinear form
    \[
        \langle\cdot,\cdot\rangle
        \colon\realnumbers^3\cartesianprod\realnumbers^3\fromto\realnumbers
    \]
    defined by
    \begin{align*}
        \langle (x_1,x_2,x_3),(y_1,y_2,y_3)\rangle
        &=
        x_1y_1+3x_2y_2+4x_3y_3\\
        &\quad
        +x_1y_2+x_2y_1-2x_2y_3-2x_3y_2.
    \end{align*}
    Decide whether it is linear in both variables, symmetric, and positive
    definite.
\end{snippetexercise}

\begin{snippetsolution}{symmetric-bilinear-form-inner-product-solution}{A symmetric bilinear form on R3}
    The map is linear in each variable because it is a sum of bilinear terms
    \(x_iy_j\). Its \matrix in the canonical \basis is
    \[
        A=
        \begin{pmatrix}
            1 & 1 & 0\\
            1 & 3 & -2\\
            0 & -2 & 4
        \end{pmatrix},
    \]
    which is symmetric. Hence the form is a \symmetricbilinearform.

    The associated \quadraticform is
    \[
        Q(x_1,x_2,x_3)
        =
        x_1^2+3x_2^2+4x_3^2+2x_1x_2-4x_2x_3.
    \]
    The \characteristicpolynomial of \(A\) is
    \[
        p_A(t)
        =
        \det(t\identmatrix{3}-A)
        =
        (t-2)(t^2-6t+2).
    \]
    The \eigenvalue[eigenvalues] are
    \[
        2,
        \qquad
        3+\sqrt{7},
        \qquad
        3-\sqrt{7}.
    \]
    They are all positive. By the eigenvalue criterion, \(Q\) is
    \positivedefiniteform[positive definite]. Therefore
    \(\langle\cdot,\cdot\rangle\) is an \innerproduct on
    \(\realnumbers^3\).
\end{snippetsolution}

\end{document}
